\PassOptionsToPackage{unicode}{hyperref}
\PassOptionsToPackage{hyphens}{url}
\documentclass[11pt,a4paper]{article}

% ── Codificação e língua ──────────────────────────────────────────────────────
\usepackage[T1]{fontenc}
\usepackage[utf8]{inputenc}
\usepackage[brazil]{babel}

% ── Tipografia ────────────────────────────────────────────────────────────────
\usepackage{lmodern}
\usepackage{microtype}

% ── Cores ─────────────────────────────────────────────────────────────────────
\usepackage{xcolor}
\definecolor{javaOrange}{HTML}{E76F00}
\definecolor{javaDeep}{HTML}{BF5700}
\definecolor{javaBlue}{HTML}{1565C0}
\definecolor{javaTeal}{HTML}{00695C}
\definecolor{javaAmber}{HTML}{B45309}
\definecolor{javaInk}{HTML}{2B2140}
\definecolor{javaCodeBg}{HTML}{FFF8F0}
\definecolor{javaRule}{HTML}{FFD7B5}
\definecolor{javaShade}{HTML}{FFF3E8}

% ── Caixas e código ───────────────────────────────────────────────────────────
\usepackage{tcolorbox}
\tcbuselibrary{skins,breakable,listings}
\usepackage{listings}
\usepackage{fvextra}

\lstdefinestyle{javastyle}{
  language=Java,
  basicstyle=\ttfamily\small,
  keywordstyle=\color{javaBlue}\bfseries,
  stringstyle=\color{javaDeep},
  commentstyle=\color{javaTeal}\itshape,
  numberstyle=\tiny\color{gray},
  breaklines=true,
  breakatwhitespace=false,
  showstringspaces=false,
  tabsize=4,
}

\newtcblisting{javacode}{
  listing only,
  listing style=javastyle,
  breakable,
  enhanced,
  colback=javaCodeBg,
  colframe=javaDeep,
  leftrule=4pt,
  rightrule=0.4pt,
  toprule=0.4pt,
  bottomrule=0.4pt,
  arc=2pt,
  fontupper=\small,
  top=4pt, bottom=4pt, left=6pt, right=6pt,
}

\newtcbox{\inlinecode}{
  on line,
  colback=javaCodeBg,
  colframe=javaRule,
  boxrule=0.5pt,
  arc=2pt,
  fontupper=\ttfamily\small\color{javaDeep},
  boxsep=1pt, left=2pt, right=2pt, top=1pt, bottom=1pt,
}

% ── Layout ────────────────────────────────────────────────────────────────────
\usepackage[margin=2.4cm, top=2.8cm, bottom=2.8cm]{geometry}

% ── Cabeçalho / rodapé ────────────────────────────────────────────────────────
\usepackage{fancyhdr}
\pagestyle{fancy}
\fancyhf{}
\renewcommand{\headrulewidth}{0.4pt}
\renewcommand{\headrule}{\hbox to\headwidth{\color{javaRule}\leaders\hrule height \headrulewidth\hfill}}
\fancyhead[L]{\small\sffamily\color{javaDeep} Java Programming MOOC}
\fancyhead[R]{\small\sffamily\color{javaDeep} Resoluções · Lição 6}
\fancyfoot[C]{\small\sffamily\color{gray}\thepage}

% ── Seções ────────────────────────────────────────────────────────────────────
\usepackage{titlesec}
\titleformat{\section}
  {\sffamily\Large\bfseries\color{javaOrange}}
  {\thesection}{0.7em}{}
  [\vspace{2pt}{\color{javaRule}\titlerule[1.2pt]}]
\titleformat{\subsection}
  {\sffamily\large\bfseries\color{javaDeep}}
  {\thesubsection}{0.6em}{}
\titleformat{\subsubsection}
  {\sffamily\normalsize\bfseries\color{javaTeal}}
  {\thesubsubsection}{0.5em}{}

% ── Listas ────────────────────────────────────────────────────────────────────
\usepackage{enumitem}
\setlist[itemize]{itemsep=2pt, topsep=4pt, parsep=0pt}
\setlist[enumerate]{itemsep=2pt, topsep=4pt, parsep=0pt}

% ── Tabelas ───────────────────────────────────────────────────────────────────
\usepackage{booktabs}
\usepackage{array}
\usepackage{colortbl}
\usepackage{longtable}

% ── Hiperlinks ────────────────────────────────────────────────────────────────
\usepackage{hyperref}
\hypersetup{
  colorlinks=true,
  linkcolor=javaBlue,
  urlcolor=javaBlue,
  citecolor=javaBlue,
  pdftitle={Programação em Java --- Resoluções --- Lição 6},
  pdfauthor={Java Programming MOOC · Universidade de Helsinque},
  pdflang={pt-BR},
}

% ── Caixa de destaque (dica/atenção) ─────────────────────────────────────────
\newtcolorbox{destaque}[1][Observação]{
  enhanced, breakable,
  colback=javaShade,
  colframe=javaOrange,
  leftrule=4pt, rightrule=0.4pt, toprule=0.4pt, bottomrule=0.4pt,
  arc=2pt,
  fonttitle=\sffamily\bfseries\color{javaOrange},
  title=#1,
  top=4pt, bottom=4pt,
}

% ── Título personalizado ──────────────────────────────────────────────────────
\makeatletter
\newcommand{\subtitle}[1]{\gdef\@subtitle{#1}}
\renewcommand{\maketitle}{%
  \begin{center}
    {\color{javaRule}\rule{\linewidth}{1pt}}\\[6pt]
    {\Huge\sffamily\bfseries\color{javaOrange} \@title}\\[6pt]
    \ifx\@subtitle\undefined\else
      {\Large\sffamily\color{javaDeep} \@subtitle}\\[4pt]
    \fi
    {\small\sffamily\color{gray} \@author}\\[2pt]
    {\small\sffamily\color{gray} \@date}\\[6pt]
    {\color{javaRule}\rule{\linewidth}{1pt}}
  \end{center}
  \vspace{1em}
}
\makeatother

% ─────────────────────────────────────────────────────────────────────────────
\title{Programação em Java — Resoluções}
\subtitle{Lição 6}
\author{Java Programming MOOC · Universidade de Helsinque — tradução para o português}
\date{2026}
% ─────────────────────────────────────────────────────────────────────────────

\begin{document}
\maketitle
\tableofcontents
\vspace{1.5em}

Os testes usam JUnit 4, do mesmo jeito que foi apresentado na lição (\inlinecode{org.junit.Test} e \inlinecode{org.junit.Assert}).

% ─────────────────────────────────────────────────────────────────────────────
\section{Exercício 1 --- Biblioteca}

\begin{javacode}
public class Livro {
    private String titulo;
    private String autor;

    public Livro(String titulo, String autor) {
        this.titulo = titulo;
        this.autor = autor;
    }

    public String getTitulo() {
        return titulo;
    }

    public String getAutor() {
        return autor;
    }

    @Override
    public String toString() {
        return titulo + " - " + autor;
    }
}
\end{javacode}

\begin{javacode}
import java.util.ArrayList;
import java.util.Collections;
import java.util.Comparator;

public class Biblioteca {
    private ArrayList<Livro> livros;

    public Biblioteca() {
        this.livros = new ArrayList<>();
    }

    public void adicionarLivro(Livro livro) {
        livros.add(livro);
    }

    public ArrayList<Livro> buscarPorAutor(String autor) {
        ArrayList<Livro> encontrados = new ArrayList<>();
        for (Livro livro : livros) {
            if (livro.getAutor().equals(autor)) {
                encontrados.add(livro);
            }
        }
        return encontrados;
    }

    public ArrayList<Livro> listarTodosOrdenados() {
        ArrayList<Livro> ordenados = new ArrayList<>(livros);
        Collections.sort(ordenados, Comparator.comparing(Livro::getTitulo));
        return ordenados;
    }
}
\end{javacode}

\begin{javacode}
import org.junit.Before;
import org.junit.Test;
import static org.junit.Assert.*;

public class BibliotecaTest {

    private Biblioteca biblioteca;

    @Before
    public void setUp() {
        biblioteca = new Biblioteca();
        biblioteca.adicionarLivro(new Livro("Torto Arado", "Itamar Vieira Junior"));
        biblioteca.adicionarLivro(new Livro("Dom Casmurro", "Machado de Assis"));
        biblioteca.adicionarLivro(new Livro("Quincas Borba", "Machado de Assis"));
    }

    @Test
    public void adicionarLivroAumentaTamanho() {
        Biblioteca vazia = new Biblioteca();
        assertEquals(0, vazia.listarTodosOrdenados().size());

        vazia.adicionarLivro(new Livro("1984", "George Orwell"));
        assertEquals(1, vazia.listarTodosOrdenados().size());
    }

    @Test
    public void buscarPorAutorRetornaTodosOsLivrosDoAutor() {
        assertEquals(2, biblioteca.buscarPorAutor("Machado de Assis").size());
        assertEquals(1, biblioteca.buscarPorAutor("Itamar Vieira Junior").size());
        assertEquals(0, biblioteca.buscarPorAutor("Autor Inexistente").size());
    }

    @Test
    public void listarTodosOrdenadosRetornaEmOrdemAlfabetica() {
        var ordenados = biblioteca.listarTodosOrdenados();

        assertEquals("Dom Casmurro", ordenados.get(0).getTitulo());
        assertEquals("Quincas Borba", ordenados.get(1).getTitulo());
        assertEquals("Torto Arado", ordenados.get(2).getTitulo());
    }
}
\end{javacode}

% ─────────────────────────────────────────────────────────────────────────────
\section{Exercício 2 --- Separação de responsabilidades}

\begin{javacode}
import java.util.ArrayList;

public class ColetorDeNumeros {
    private ArrayList<Integer> numeros;

    public ColetorDeNumeros() {
        this.numeros = new ArrayList<>();
    }

    public void adicionar(int numero) {
        numeros.add(numero);
    }

    public int soma() {
        int soma = 0;
        for (int n : numeros) {
            soma += n;
        }
        return soma;
    }

    public double media() {
        if (numeros.isEmpty()) {
            return 0;
        }
        return (double) soma() / numeros.size();
    }
}
\end{javacode}

\begin{javacode}
import java.util.Scanner;

public class InterfaceColetorDeNumeros {
    public static void main(String[] args) {
        Scanner scanner = new Scanner(System.in);
        ColetorDeNumeros coletor = new ColetorDeNumeros();

        while (true) {
            String linha = scanner.nextLine();
            if (linha.equals("fim")) {
                break;
            }
            coletor.adicionar(Integer.valueOf(linha));
        }

        System.out.println("Soma: " + coletor.soma());
        System.out.println("Média: " + coletor.media());
    }
}
\end{javacode}

\begin{destaque}
Com essa separação, \inlinecode{ColetorDeNumeros} pode ser testada com JUnit sem precisar simular entrada de usuário --- toda a lógica de soma e média fica isolada da leitura via \inlinecode{Scanner}.
\end{destaque}

% ─────────────────────────────────────────────────────────────────────────────
\section{Exercício 3 --- TDD: Conta bancária}

Seguindo TDD, os testes vêm primeiro:

\begin{javacode}
import org.junit.Before;
import org.junit.Test;
import static org.junit.Assert.*;

public class ContaBancariaTest {

    private ContaBancaria conta;

    @Before
    public void setUp() {
        conta = new ContaBancaria();
    }

    @Test
    public void depositarAumentaSaldo() {
        conta.depositar(100);
        assertEquals(100, conta.getSaldo(), 0.001);
    }

    @Test
    public void sacarComSaldoSuficienteDiminuiSaldoERetornaTrue() {
        conta.depositar(100);

        boolean resultado = conta.sacar(40);

        assertTrue(resultado);
        assertEquals(60, conta.getSaldo(), 0.001);
    }

    @Test
    public void sacarSemSaldoSuficienteNaoAlteraSaldoERetornaFalse() {
        conta.depositar(50);

        boolean resultado = conta.sacar(100);

        assertFalse(resultado);
        assertEquals(50, conta.getSaldo(), 0.001);
    }

    @Test
    public void saldoNuncaFicaNegativo() {
        conta.sacar(500); // conta começa com saldo 0
        assertTrue(conta.getSaldo() >= 0);
    }
}
\end{javacode}

E só então a implementação, feita para satisfazer os testes acima:

\begin{javacode}
public class ContaBancaria {
    private double saldo;

    public ContaBancaria() {
        this.saldo = 0;
    }

    public void depositar(double valor) {
        saldo += valor;
    }

    public boolean sacar(double valor) {
        if (valor > saldo) {
            return false;
        }
        saldo -= valor;
        return true;
    }

    public double getSaldo() {
        return saldo;
    }
}
\end{javacode}

% ─────────────────────────────────────────────────────────────────────────────
\section{Exercício 4 --- Programa complexo: inventário}

\begin{javacode}
public class Item {
    private String nome;
    private int quantidade;
    private double precoUnitario;

    public Item(String nome, int quantidade, double precoUnitario) {
        this.nome = nome;
        this.quantidade = quantidade;
        this.precoUnitario = precoUnitario;
    }

    public String getNome() {
        return nome;
    }

    public int getQuantidade() {
        return quantidade;
    }

    public void setQuantidade(int quantidade) {
        this.quantidade = quantidade;
    }

    public double getValorTotal() {
        return quantidade * precoUnitario;
    }

    @Override
    public String toString() {
        return nome + " (qtd: " + quantidade + ", unit: R$" + precoUnitario + ")";
    }
}
\end{javacode}

\begin{javacode}
import java.util.ArrayList;

public class Inventario {
    private ArrayList<Item> itens;

    public Inventario() {
        this.itens = new ArrayList<>();
    }

    public void adicionar(Item item) {
        itens.add(item);
    }

    public boolean remover(String nome) {
        Item item = buscar(nome);
        if (item == null) {
            return false;
        }
        itens.remove(item);
        return true;
    }

    public Item buscar(String nome) {
        for (Item item : itens) {
            if (item.getNome().equals(nome)) {
                return item;
            }
        }
        return null;
    }

    public double valorTotal() {
        double total = 0;
        for (Item item : itens) {
            total += item.getValorTotal();
        }
        return total;
    }

    public ArrayList<Item> listar() {
        return itens;
    }
}
\end{javacode}

\begin{javacode}
import org.junit.Before;
import org.junit.Test;
import static org.junit.Assert.*;

public class InventarioTest {

    private Inventario inventario;

    @Before
    public void setUp() {
        inventario = new Inventario();
        inventario.adicionar(new Item("Parafuso", 100, 0.10));
        inventario.adicionar(new Item("Martelo", 5, 25.00));
    }

    @Test
    public void buscarEncontraItemExistente() {
        assertNotNull(inventario.buscar("Martelo"));
        assertNull(inventario.buscar("Chave de fenda"));
    }

    @Test
    public void removerRetornaTrueQuandoItemExiste() {
        assertTrue(inventario.remover("Parafuso"));
        assertNull(inventario.buscar("Parafuso"));
    }

    @Test
    public void removerRetornaFalseQuandoItemNaoExiste() {
        assertFalse(inventario.remover("Item Fantasma"));
    }

    @Test
    public void valorTotalSomaCorretamenteTodosOsItens() {
        // 100 * 0.10 + 5 * 25.00 = 10 + 125 = 135
        assertEquals(135.0, inventario.valorTotal(), 0.001);
    }
}
\end{javacode}

\begin{javacode}
import java.util.Scanner;

public class InterfaceInventario {
    public static void main(String[] args) {
        Scanner scanner = new Scanner(System.in);
        Inventario inventario = new Inventario();

        while (true) {
            System.out.println("\n1-Adicionar 2-Remover 3-Buscar 4-Listar 5-Total 6-Sair");
            String opcao = scanner.nextLine();

            switch (opcao) {
                case "1":
                    System.out.print("Nome: ");
                    String nome = scanner.nextLine();
                    System.out.print("Quantidade: ");
                    int quantidade = Integer.valueOf(scanner.nextLine());
                    System.out.print("Preço unitário: ");
                    double preco = Double.valueOf(scanner.nextLine());
                    inventario.adicionar(new Item(nome, quantidade, preco));
                    break;
                case "2":
                    System.out.print("Nome do item a remover: ");
                    boolean removido = inventario.remover(scanner.nextLine());
                    System.out.println(removido ? "Removido." : "Item não encontrado.");
                    break;
                case "3":
                    System.out.print("Nome do item a buscar: ");
                    Item encontrado = inventario.buscar(scanner.nextLine());
                    System.out.println(encontrado != null ? encontrado : "Item não encontrado.");
                    break;
                case "4":
                    for (Item item : inventario.listar()) {
                        System.out.println(item);
                    }
                    break;
                case "5":
                    System.out.println("Valor total: R$" + inventario.valorTotal());
                    break;
                case "6":
                    return;
                default:
                    System.out.println("Opção inválida.");
            }
        }
    }
}
\end{javacode}

\end{document}
